% ============================================================
%  Springer LNCS Format Paper
%  Title: Development and Performance Validation of an Outdoor
%         Autonomous Navigation Framework for Unmanned Ground Vehicles
%  Authors: Team 13, KLE Technological University
% ============================================================

\documentclass[runningheads]{llncs}

% ===== PACKAGES =====
\usepackage{amsmath,amssymb,amsfonts}
\usepackage{algorithmic}
\usepackage{graphicx}
\usepackage{textcomp}
\usepackage{xcolor}
\usepackage{url}
\usepackage{booktabs}
\usepackage{multirow}
\usepackage{array}
\usepackage{float}
\usepackage{subcaption}
\usepackage{hyperref}
\usepackage{tikz}
\usetikzlibrary{shapes.geometric, arrows.meta, positioning, calc, fit}

\hypersetup{
    colorlinks=true,
    linkcolor=black,
    citecolor=black,
    urlcolor=blue
}

% Correct bad hyphenation here
\hyphenation{op-tical net-works semi-conduc-tor}

\begin{document}

% ===== TITLE =====
\title{Development and Performance Validation of an Outdoor Autonomous Navigation Framework for Unmanned Ground Vehicles}
\titlerunning{Outdoor Autonomous Navigation Framework for UGVs}

% ===== AUTHORS =====
\author{Shrigovinda T. Kulkarni\inst{1} \and
Tanmay Hiremath\inst{1} \and
Nisar Ahmed\inst{1} \and
Bhoomika Hotagi\inst{1} \and
Rakesh P Tapaskar\inst{1} \and
Shreyansh Hajare\inst{1} \and
Manikant Binkadakatti\inst{1}}

\authorrunning{S.T. Kulkarni et al.}

\institute{Department of Automation and Robotics Engineering, CARR LAB, KLE Technological University,
Hubballi 580031, Karnataka, India\\
\email{\{shrigovindak, tanmayhiremath, p.nisar6154, bhoomikahotagi061\}@gmail.com, rptapaskar@kletech.ac.in, shreyanshhajare@gmail.com, manikantbinkadakatti@gmail.com}}

\maketitle

% ===== ABSTRACT =====
\begin{abstract}
Outdoor surveillance and monitoring tasks can be really tough and time consuming for people to do. They can also be unsafe for the people doing them. The robots that already exist to help with these tasks are often very expensive. Need someone to watch them all the time. These robots can also have a time working well in different outdoor environments.

This paper is about the design and development of an Unmanned Ground Vehicle or UGV for short that can navigate and do surveillance outside on its own. The UGV we made uses a Raspberry Pi 5 controller. Has many sensors, like an RPLiDAR to detect obstacles a GPS module to figure out where it is an Inertial Measurement Unit to estimate its orientation and a depth camera to see its surroundings. The UGV is powered by a Li-ion battery. Has four Brushless DC hub motors that are controlled by RKI 9206 motor drivers. We also set up a WebSocket-based telemetry server so we can monitor the UGVs status and how it is working in time from a distance.

The system we made has four parts: sensing, perception, planning and control. We tested the UGV on outdoor surfaces, like grass and gravel. The results showed that the UGV can navigate on its own between set points detect obstacles and work stably. The UGV can move at speeds from 0.5 meters per second to 1.5 meters per second. Can work for over 30 minutes on one charge.

The Unmanned Ground Vehicle we developed is a cost scalable solution for outdoor autonomous navigation and surveillance tasks. In the future we plan to make the UGV even better, by adding Simultaneous Localization and Mapping techniques and reinforcement learning-based path planning to improve how it navigates and adapts to environments. The Unmanned Ground Vehicle will be able to navigate better and work well in many different situations.

\keywords{Autonomous Navigation \and BLDC Hub Motor \and GPS--IMU Sensor Fusion \and LiDAR \and Obstacle Avoidance \and Path Planning \and Raspberry Pi \and Remote Monitoring \and Unmanned Ground Vehicle}
\end{abstract}

% ================================================================
%                      I. INTRODUCTION
% ================================================================
\section{Introduction}

The use of robotic systems in outdoor environments has become really popular in the last few years. This is because they can be used in different areas such as agriculture, surveillance, disaster response and defense. Unmanned Ground Vehicles are an option for monitoring places that are dangerous or hard for people to reach.

Even though autonomous vehicle technology has come a way there are still some big challenges to overcome. Most of the systems are very expensive and use things like multi-beam LiDAR, high-precision RTK-GPS and special computing platforms. This makes them too expensive for universities and small groups to use. Also many of these systems only work well on roads and have trouble with grass, gravel and hills.

The main problems with Unmanned Ground Vehicles navigating outdoors are safety, cost adapting to terrain and communication. For example Unmanned Ground Vehicles are used in environments which can put people in danger. Unmanned Ground Vehicles are also very expensive which makes them hard to use in situations. Unmanned Ground Vehicles have trouble with terrain, like grass and hills.. Unmanned Ground Vehicles can have communication problems when trying to control them from far away.

This paper talks about an Unmanned Ground Vehicle platform that is low-cost and has a special navigation system. The Unmanned Ground Vehicle platform has five features. First it uses a design method to make sure all the parts work well together. Second it has a four-layer software system that includes sensing, perception, planning and control. Third it uses GPS and other sensors to detect obstacles. Fourth it has a real-time telemetry system that lets people control it from away. Fifth it has been tested on different types of outdoor terrain and can navigate to specific points on its own.

The rest of this paper is organized into sections. Section II looks at work that has been done on Unmanned Ground Vehicle navigation. Section III explains the design, hardware and software of the Unmanned Ground Vehicle platform. Section IV talks about the results of the tests and how well the Unmanned Ground Vehicle platform worked. Section V summarizes the paper.

% ================================================================
%                   II. LITERATURE REVIEW
% ================================================================
\section{Literature Review}
\label{sec:literature}

The field of self-driving UGV navigation has moved away from using one sensor. Now it uses sensors and systems working together.Here are the main points we found in our research.They helped us design our system.We looked at studies to understand what works and what doesn't in UGV navigation.These studies showed us that using sensors is better than using just one.Our system uses this information to navigate UGV navigation is a task.It requires sensors and systems.The design of our system is based on what we learned from these studies.We used sensors to make sure our UGV can navigate safely.This is important for UGV navigation.UGV navigation is an area of research.It has applications, in the real world.Our system is designed to help with these applications.

\subsection{Autonomous Vehicle Architectures}

Software for systems needs to be organized in a good way.A simple system with four layers is often used: it can sense things it can understand what it senses it can plan what to do. It can control itself.This system uses equipment like Velodyne LiDAR and GNSS+IMU to know where it is.This system works well but it is expensive and needs maps of the area so it is not good for people who do not have a lot of money or for schools.To make it cheaper people have tried to make the system more simple so it can be. Used more easily.They have also tried to make the system work better with power so it can be used in more places.If we use very cheap parts, like Arduino and simple sensors the system is not powerful enough to understand things that are happening and make plans when it is moving around outside.Autonomous systems, like these need a system that is not too expensive.
Autonomous systems need to be able to sense and understand things and make plans so they can work well in the world.

\subsection{Multi-Sensor Fusion and Localization}

Reliable localization in unstructured environments requires fusing complementary sensory inputs. Research has explored integrating RTK-GPS, IMU, and LiDAR using edge computing nodes to reduce latency~\cite{seelam2025iot}, and employing non-linear sliding-mode controllers to mitigate terrain-induced slip~\cite{gonzalez2025robust}. Multi-modal sensor fusion with memory-guided networks has also shown high real-time tracking performance in dynamic environments~\cite{han2025real}. Under extreme conditions where ground features are obscured, cooperative air-ground frameworks leveraging UWB trilateration~\cite{yue2025uav} or visual-SLAM and terrain segmentation models~\cite{gupta2025drone} achieve high localization and mapping accuracy, though at the cost of high system complexity.

\subsection{Advanced Path Planning and Obstacle Avoidance}

Safe navigation requires transitioning from reactive obstacle avoidance to global, risk-aware path planning. While low-cost acoustic-based reactive avoidance is suitable for slow rovers~\cite{desimone2018obstacle}, fast-moving outdoor UGVs demand active planners. Fusing Artificial Potential Fields (APF) with A* search algorithms provides risk-aware navigation~\cite{jiang2024risk}, while tactile probes paired with model predictive control allow crop-canopy navigation where visual and GPS signals are attenuated~\cite{gronewold2025tactile}. For unstructured terrain, deep reinforcement learning (DRL) using curriculum training has been applied to learn traversability mapping from 3D LiDAR~\cite{sanchez2023reinforcement} and follow forest trails~\cite{tibermacine2026autonomous}. Additionally, integrating visual-inertial odometry with real-time traversability estimation enables tactical motion planning over steep, uneven gradients~\cite{mukherjee2025terrain}.

\subsection{UGV Subsystem Design and Applications}

Present-day UGVs employ specific active subsystems together with their motion subsystem. Mathematical modeling and kinematic verification were used for torque and stability assessment of the chassis climbing ability.~\cite{dolya2025development}. Active Active load carrying devices, such as docking systems for charging automation.
~\cite{jung2026autonomous}, active cable deployment systems for mobile microgrids.~\cite{naglak2021cable},and communication node ejection systems for keeping an ad-hoc network operational during search and rescue operations~\cite{fisher2025developing}, Illustrate how multifunctional increase in the current field  of UGV platforms.

\subsection{Research Gaps and Motivation}

However, there still exist several critical areas that have not been sufficiently addressed yet: (i) high-quality navigation systems are too expensive for being used in academic researches; (ii) inexpensive GPS/IMU/LiDAR-based systems fail due to noises and wheel slipping under strict computing constraints; and (iii) remote telemetry is not taken into account enough, resulting in high latency of packets. These issues are solved in this paper through creating affordable UGV with COTS components combined with a highly-efficient software design.

% ================================================================
%                   III. METHODOLOGY
% ================================================================
\section{Methodology}
\label{sec:methodology}

This section provides the entire  methodology that covers hardware development, sensor fusion design, and software pipeline Deployment. The UGV system design follows an engineering methodology involving the process of sensor fusion design, hardware development, and software pipeline development.

\subsection{System Block Diagram}

The UGV design is created to provide a low-cost, strong, and modular system platform that can operate in unstructured outdoor environments. The architecture of the system is broken down into separate functional systems in order to guarantee proper isolation of the computation process, sensor information, power management, and physical actuators.

 Fig.~\ref{fig:block_diagram} shows the block diagram of the proposed UGV architecture. It clearly depicts the hierarchical connections between the processing unit (Raspberry Pi~5), the sensors unit (GPS, IMU, LiDAR, depth camera), the power unit (battery, boost/buck converters), the actuation unit (BLDC hub motors, RKI~9206 drivers) and the communications unit (Wi-Fi telemetry link to ground control station). The presented block diagram shows the signal chain, which starts with sensing and ends with actuation control, being a basis of modular four-layer software architecture.

\begin{figure}[htbp]
\centering
\includegraphics[width=0.7\columnwidth]{Block Diagram.png}
\caption{Block diagram of the proposed UGV system architecture showing sensor, processing, actuation, and communication subsystems.}
\label{fig:block_diagram}
\end{figure}

\subsection{GPS--IMU Sensor Fusion}

Localization for outdoor applications that enable autonomous waypoints navigation must be precise. Single stand-alone GPS units experience problems with multi-pathing errors, signal loss due to foliage, as well as low update rate (about 1--5\,Hz). Meanwhile, IMU dead-reckoning systems develop drift due to bias in accelerometers and gyros. In order to compensate for these shortcomings, a loosely coupled GPS+IMU solution is suggested.

The fusion architecture operates as follows.The MPU6050 IMU supplies orientation measurements (roll, pitch, yaw) at high frequency (50\,Hz) along with short-term positioning by double integration of accelerometer readings. The GPS L76X gives absolute positions at 1\,Hz frequency with accuracy of $\leq\pm$2\,m under open sky conditions. Complementary filter fuses both sensors' outputs; specifically, at each cycle of control loop, heading and displacement calculated by IMU are used for rapid navigation updates whereas GPS positions are utilized as correction signals to restrict IMU's drift. The fused heading estimate $\psi_{\text{fused}}$ is computed as:

\begin{equation}
\psi_{\text{fused}} = \alpha \cdot \psi_{\text{IMU}} + (1 - \alpha) \cdot \psi_{\text{GPS}}
\end{equation}

Where $\alpha$ $\epsilon$ $[0, 1]$ as an adjustable weighting factor (chosen as 0.85 in this implementation), which gives more weightage to IMU in the short run but depends on GPS in the long run for accuracy. Fused positions can also be modeled in a similar correction fashion, where displacements based on the IMU are corrected by GPS positions in order to avoid drifting indefinitely.

Such a sensor fusion technique makes sure of heading and positioning estimations even in short periods of GPS outage, making navigation on outdoor terrains possible.

\subsection{Hardware Platform}

The architecture of UGV hardware platform consists of a combination of mechanical, electrical, and sensors systems to form an effective mobile platform. Figure~\ref{fig:system_arch} demonstrates  the structure of the overall system framework design, which includes the autonomous navigation pipeline consisting of four layers, the hardware platform, and the ground control station.

\begin{figure}[htbp]
\centering
\includegraphics[width=0.7\columnwidth]{UGV Framework.png}
\caption{System framework design of the proposed UGV autonomous navigation platform.}
\label{fig:system_arch}
\end{figure}

\subsubsection{Mechanical Subsystem.}

The chassis is made from steel. It is 400 mm x 600 mm x 20 mm in size. This gives the chassis strength for use on terrain. The acrylic sheet is 400 mm x 600 mm x 5 mm, in size. It is used to hold the electronics in place. The assembly uses M12 x 50 mm fasteners to keep the components in place. The M12 x 50 mm fasteners are used for mounting the components to the chassis and the acrylic sheet.

 

\subsubsection{Drive System.}

The locomotion  assembly unit has 4 BLDC hub motors of 350\,W capacity for all-wheel drive. Each motor operates using an independent RKI~9206 motor controller that works within 31--36\,V power input range, able to supply up to 16\,A per motor during maximum load. The two wheel differential drive design (rear wheels active and front wheels passive) was chosen due to its simplicity.

\subsubsection{Power Subsystem.}

A TATTU 30000\textbackslash{},mAh Li-ion battery is the power source. It provides 24\textbackslash{},V voltage.

The power distribution architecture has parts:

Boost Converters (4$\times$): Four boost converters are used to increase the battery voltage from 22--24\,V to 36\,V for the motor driver input. The output voltage is regulated within a $\pm$5\% tolerance.

Buck Converter: This is a DC-DC 5\textbackslash{},V 5\textbackslash{},A regulator. It gives logic-level power to the Raspberry Pi\~5, sensors and communication modules.

4-Channel Relay Module: This module helps control the direction of the motor drivers. It switches control signals under software control.

 


\subsubsection{Sensor Suite.}

The multi-sensor perception system is made up of these things:

1. RP LiDAR: This is a 360 degree 2D laser scanner. The ``RP LiDAR'' can detect things that're between 0.2 and 10 meters away. We use the ``RP LiDAR'' for finding obstacles and making a map of the environment.

2. GPS L76X: The  GPS L76X  is a GPS module. It tells us where we are on the earth. The GPS L76X is usually right within 2 meters.

3. MPU6050 IMU: The MPU6050 IMU tells us which way we are facing. It also tells us how fast we are turning and how fast we are moving. The  MPU6050 IMU  sends us this information at 50 times per second.

4. Depth Camera D435i: The  Depth Camera D435i is a camera. It is an Intel Real Sense D435i depth camera. The  Depth Camera D435i  helps us see what is around us, by sending us pictures and videos. We use the  Depth Camera D435i  to know what is going on.

 

\subsubsection{Computing and Communication.}

The Raspberry Pi~5 (16\,GB RAM) is used for central computation purposes, including the navigation algorithm, sensor drivers, and telemetry server. The wireless connection is achieved using Wi-Fi that ensures the telemetry transfer in real-time with a delay of less than 5 seconds.

\subsection{Software Architecture}

The design of the software architecture is based on a modular four-layer pipeline which uses the Robot Operating System 2 (ROS2) Humble Hawksbill framework on the Raspberry Pi~5 platform. The software architecture design features an extensible nature, clear separation of the nodes, and real-time performance with a control loop frequency of 10\,Hz.

Modularity is achieved through compiling sensors drivers, perception, and planning models into different ROS2 nodes.Data transfer between these nodes relies on standard ROS2 publisher-subscriber communication channels. The primary ROS2 topics utilized in the system are:
\begin{itemize}\setlength{\itemsep}{0.2ex}\parskip0pt\parsep0pt
    \item \texttt{/scan} (\textit{sensor\_msgs/msg/LaserScan}): Range and angle data published by the RPLidar driver node at 10\,Hz.
    \item \texttt{/gps/fix} (\textit{sensor\_msgs/msg/NavSatFix}): Raw geolocation parameters (latitude, longitude, altitude) published by the GPS driver node at 1\,Hz.
    \item \texttt{/imu/data\_raw} (\textit{sensor\_msgs/msg/Imu}): Raw accelerometer and gyroscope measurements published by the IMU driver node at 50\,Hz.
    \item \texttt{/camera/depth/image\_rect\_raw} (\textit{sensor\_msgs/msg/Image}): Rectified depth stream published by the D435i camera driver node.
    \item \texttt{/odom/fused} (\textit{nav\_msgs/msg/Odometry}): Fused pose and orientation updates published by the sensor fusion perception node.
    \item \texttt{/cmd\_vel} (\textit{geometry\_msgs/msg/Twist}): Target linear and angular velocity commands published by the path planning node.
    \item \texttt{/motor\_cmds} (\textit{std\_msgs/msg/Float32MultiArray}): Duty cycle commands published by the control node to the RKI~9206 drivers.
\end{itemize}

\begin{figure}[htbp]
\centering
\includegraphics[width=0.42\columnwidth]{Flowchart.png}
\caption{Operational flowchart and navigation decision pipeline of the UGV system.}
\label{fig:flowchart}
\end{figure}

\subsubsection{Sensing Layer.}

The sensing layer interfaces with all hardware sensors through dedicated driver modules:

\begin{itemize}\setlength{\itemsep}{0.2ex}\parskip0pt\parsep0pt
    \item \textbf{RPLidar Driver}: Communicates with the RP~LiDAR via USB serial (\texttt{/dev/ttyUSB0}), collecting 360\textdegree\ scan data at 2\textdegree\ angular resolution, yielding 180 range measurements per scan.
    \item \textbf{GPS Driver}: Parses NMEA sentences (\texttt{\$GPRMC}, \texttt{\$GPGGA}) from the GPS module via serial interface at 9600~baud rate.
    \item \textbf{IMU Driver}: Reads orientation quaternions and acceleration vectors at 50\,Hz via I\textsuperscript{2}C interface.
    \item \textbf{Camera Driver}: Captures RGB and depth frames using OpenCV and Intel RealSense SDK (pyrealsense2).
\end{itemize}

All sensor modules implement graceful fallback to simulation mode when hardware is unavailable, enabling software development and testing without physical sensors.

\subsubsection{Perception Layer.}

The perception layer processes raw sensor data to extract meaningful environmental information. The minimum distance to obstacles is calculated within a 40-degree field of view centered on the front sector of the LiDAR, and an obstacle is flagged as ``NEAR'' when this distance falls below a threshold of 1.0\,m. LiDAR-based obstacle detection employs ray-circle intersection testing to calculate the distance and collision parameters between the robot and detected obstacles, facilitating dynamic hazard assessment.

\subsubsection{Planning Layer.}

The planning layer implements waypoint-based goal navigation with reactive obstacle avoidance, drawing insights from modern off-road path planners~\cite{jiang2024risk, mukherjee2025terrain}. Given a goal position and the current robot position, the target heading is computed using the arctangent of the relative coordinate offsets. Smooth heading control is achieved by calculating and clamping the angular correction to prevent rapid steering oscillations. Goal arrival is detected when the Euclidean distance to the target is less than 0.2\,m.

\subsubsection{Control Layer.}

The control layer translates planned motion commands into motor driver signals. The system supports two control modes:

\begin{enumerate}\setlength{\itemsep}{0.2ex}\parskip0pt\parsep0pt
    \item \textbf{Autonomous Mode}: The planning layer generates linear velocity and angular corrections, which are mapped to differential wheel speeds for the BLDC hub motors through the RKI~9206 drivers.
    \item \textbf{Teleoperation Mode}: The operator sends drive commands (linear, angular) from the ground control station via WebSocket, directly controlling motor speeds.
\end{enumerate}

\subsubsection{Telemetry and Communication.}

A WebSocket server implements the real-time telemetry system, broadcasting sensor data and camera feeds at 10\,Hz to connected ground control clients. In accordance with robust telemetry protocols established for remote outdoor networks~\cite{fisher2025developing, memet2025adaptive},Telemetry includes:

\begin{itemize}\setlength{\itemsep}{0.2ex}\parskip0pt\parsep0pt
    \item Robot coordinates $(x, y)$ and heading $\psi$
    \item GPS coordinates ( latitude, longitude, satellites count)
    \item LiDAR range data (180 points)
    \item Depth sensor data (distance, NEAR/FAR classification)
    \item Navigation path waypoints 
\end{itemize}

Server can handle several clients simultaneously and uses asynchronous I/O via Python \texttt{asyncio}.

% ================================================================
%              VI. EXPERIMENTAL RESULTS
% ================================================================
\section{Experimental Results and Validation}
\label{sec:results}

The UGV platform was tested in outdoor experiments. These experiments took place on the KLE Technological University campus. The experiments had types of ground. There were areas, grass, gravel paths and slopes. The UGV platform worked well on all these surfaces. The campus had everything needed for the tests. UGV platform testing is important.The KLE Technological University campus was ideal, for testing the UGV platform.

 

\subsection{Localization Accuracy}

The GPS–IMU system was assessed through commanding the UGV to move between specified waypoints and assessing the accuracy of the GPS path obtained from the actual path. The accuracy achieved by the system when the GPS condition is good and the visible satellites are 12 is $\pm$2\,m.

\subsection{Obstacle Detection Performance}

The RP~LiDAR successfully detected obstacles within the 0.2--10\,m range over the entire 360\textdegree\ viewing angle. It was evaluated for the ability to detect static obstacles of various dimensions situated between 0.3\,m and 8\,m away. It successfully detected more than 50\% of the obstacles according to the design requirement (FR-13). Better accuracy was achieved for obstacles that were larger than 0.3\,m in diameter located up to 5\,m away. The performance metrics and the accuracies of the components are tabulated in Table~\ref{tab:component_accuracy}.

\begin{table}[htbp]
\caption{Component Performance and Measurement Accuracy}
\label{tab:component_accuracy}
\centering
\resizebox{\columnwidth}{!}{%
\begin{tabular}{llll}
\toprule
\textbf{Component} & \textbf{Measurement Parameter} & \textbf{Operational Range} & \textbf{Measured Accuracy} \\
\midrule
GPS L76X & Global Position & Global (open-sky) & $\pm$2.0\,m (CEP) \\
IMU MPU6050 & Heading \& Angular Velocity & $\pm$250\textdegree/s & $\pm$1.0\textdegree\ (Fused) \\
RP LiDAR & Obstacle Distance & 0.2--10\,m & $\pm$1.0\% of range \\
Depth Camera D435i & Frontal Obstacle Depth & 0.3--10\,m & $<$2.0\% at 2\,m \\
BLDC Motor Drivers & Output Voltage Regulation & 31--36\,V & $\pm$5.0\% tolerance \\
\bottomrule
\end{tabular}%
}
\end{table}

\subsection{Terrain Adaptability}

In order to test the performance of the UGV, extensive experiments have been conducted in different outdoor environments. These experiments were carried out using the following major types of terrains:
\begin{itemize}\setlength{\itemsep}{0.2ex}\parskip0pt\parsep0pt
    \item \textbf{Flat ground and paved terrains}: This is the base terrain. Experiments conducted on flat ground and paved terrains such as concrete floor and tiled pathways revealed stable locomotion. Heading control was performed with no wheel slip, thus reaching the target speed of 0.5--1.5\,m/s. For flat ground terrains, the cross track error for waypoint tracking was less than $\pm$0.5\,m.
\item \textbf{Grass and Gravel} : Tested to test off-road durability. The robot kept navigation steady with slight reduction in speed (0.8 m/s) due to greater rolling friction and slip of gravel, which is successfully controlled by the sensor fusion feedback system.
    \item \textbf{Sloping surface} : Tested for slopes up to 15\textdegree. The robot had sufficient traction and torque distribution without motor overheating and reverse drifting.
\end{itemize}

The BLDC hub motors (4 $\times$ 350\,W) provided sufficient torque for all tested terrain conditions at the specified operating voltages of 31--36\,V.

\subsection{Communication and Telemetry}

The WebSocket-based telemetry system achieved real-time data streaming with latency below 5 seconds between the UGV and ground control station over Wi-Fi. The system supported simultaneous streaming of:
\begin{itemize}\setlength{\itemsep}{0.2ex}\parskip0pt\parsep0pt
    \item Telemetry data at 10\,Hz update rate
    \item RGB camera feed at 720p, 15\,FPS
    \item Depth camera visualization
    \item LiDAR scan data (180 points per frame)
\end{itemize}

\subsection{Battery Endurance}

The TATTU 30000\,mAh Li-ion battery provided a minimum of 30 minutes continuous operation with all motors and sensors active during autonomous navigation, meeting the design specification. The power distribution through boost converters maintained stable 36\,V output within $\pm$5\% tolerance throughout the operational window.

\subsection{Comparative Analysis}

To contextualize the performance of the developed UGV, Table~\ref{tab:literature_comparison} presents a comparative evaluation against similar systems described in existing literature across multiple operational and hardware dimensions.

\begin{table}[htbp]
\caption{Performance Comparison with Existing Literature}
\label{tab:literature_comparison}
\centering
\resizebox{\columnwidth}{!}{%
\begin{tabular}{llllll}
\toprule
\textbf{Reference} & \textbf{Loc. Acc.} & \textbf{3D Vision} & \textbf{Telemetry Link} & \textbf{Cost} & \textbf{Overall Utility} \\
\midrule
Azam et al.~\cite{azam2020system} & $<$0.1\,m & Yes & Standard RF & High & Over-engineered \\
Fernandez et al.~\cite{fernandez2019unmanned} & $\pm$5.0\,m & No & None & V. Low & Under-powered \\
Dolya et al.~\cite{dolya2025development} & $\pm$1.5\,m & No & High-Latency HTTP & Med. & Moderate (No IMU) \\
\textbf{Proposed Work} & \textbf{$\pm$2.0\,m} & \textbf{Yes} & \textbf{WebSocket (Real-time)} & \textbf{Low} & \textbf{Excellent (Optimized)} \\
\bottomrule
\end{tabular}%
}
\end{table}

% ================================================================
%                  VII. CONCLUSION
% ================================================================
\section{Conclusion}
\label{sec:conclusion}

This paper presented the development and performance validation of a low-cost outdoor autonomous navigation framework for Unmanned Ground Vehicles. The systematic engineering design methodology and GPS--IMU sensor fusion yielded an optimal platform configuration integrating a Raspberry Pi~5 with fused GPS--IMU localization, RP~LiDAR obstacle detection, BLDC hub motor drive, and WebSocket-based remote monitoring.

Experimental validation confirmed that the proposed system achieves reliable autonomous waypoint navigation on diverse outdoor terrain with GPS accuracy within $\pm$2\,m, obstacle detection within 0.2--10\,m, and operational endurance exceeding 30~minutes.

% ===== REFERENCES =====
\begin{thebibliography}{00}

\bibitem{azam2020system}
S.~Azam, F.~Munir, A.~M.~Sheri, J.~Kim, and M.~Jeon, ``System, design and experimental validation of autonomous vehicle in an unconstrained environment,'' \textit{Sensors}, vol.~20, no.~21, p.~5999, 2020.

\bibitem{fernandez2019unmanned}
S.~G.~Fernandez \textit{et~al.}, ``Unmanned and autonomous ground vehicle,'' \textit{International Journal of Electrical and Computer Engineering (IJECE)}, vol.~9, no.~5, pp.~4466--4472, 2019.

\bibitem{dolya2025development}
A.~Dolya, B.~Kolumbetov, and A.~Kalkabek, ``Development of a model and performance analysis of the characteristics of a multifunctional autonomous unmanned ground vehicle,'' \textit{Advances in Materials Science and Engineering}, vol.~2025, Art. no. 6614303, 2025.

\bibitem{jiang2024risk}
J.~Jiang, Z.~Hu, Z.~Xie, C.~Hao, \textit{et~al.}, ``A risk-aware planning framework of UGVs in off-road environment,'' \textit{IEEE Transactions on Vehicular Technology}, vol.~74, no.~3, pp.~3870--3884, 2024.

\bibitem{gadekar2023recent}
A.~Gadekar, K.~Kataria, J.~Aher, P.~Deshmukh, S.~Fulsundar, S.~Barve, and V.~Patel, ``Recent developments in modular unmanned ground vehicles: A review,'' \textit{Asia-Pacific Journal of Science and Technology}, vol.~29, no.~1, pp.~1--11, 2023.

\bibitem{khriji2016integration}
L.~Khriji \textit{et~al.}, ``An integration framework for UGV outdoor navigation system based on multi-sensor fusion,'' in \textit{Proc. International Conference on Control, Decision and Information Technologies (CoDIT)}, 2016, pp.~245--250.

\bibitem{seelam2025iot}
S.~K.~Seelam, S.~N.~Bhavanam, V.~Midasala, and E.~S.~Reddy, ``Design and Development of IoT Smart Driving Protocol for Unmanned Ground Vehicles,'' in \textit{Proc. 6th International Conference on Mobile Computing and Sustainable Informatics (ICMCSI)}, 2025, pp.~336--341.

\bibitem{gonzalez2025robust}
I.~Gonz\'{a}lez-Hern\'{a}ndez, J.~Flores, S.~Salazar, and R.~Lozano, ``Robust and Precise Navigation and Obstacle Avoidance for Unmanned Ground Vehicle,'' \textit{Sensors}, vol.~25, no.~14, p.~4334, 2025.

\bibitem{han2025real}
Z.~Han, P.~Chen, B.~Zhou, and G.~Yu, ``Real-time navigation of unmanned ground vehicles in complex terrains with enhanced perception and memory-guided strategies,'' \textit{IEEE Transactions on Vehicular Technology}, vol.~74, no.~3, pp.~3723--3735, 2025.

\bibitem{yue2025uav}
P.~Yue, J.~Xin, Y.~Huang, J.~Zhao, C.~Zhang, W.~Chen, and M.~Shan, ``UAV autonomous navigation system based on air--ground collaboration in GPS-denied environments,'' \textit{Drones}, vol.~9, no.~6, p.~442, 2025.

\bibitem{gupta2025drone}
S.~Gupta, P.~Agrawal, and P.~Gupta, ``Drone-assisted UAV-UGV collaboration for autonomous navigation in snow-covered terrain,'' \textit{HAL Preprint}, 2025.

\bibitem{desimone2018obstacle}
M.~C.~De~Simone, Z.~B.~Rivera, and D.~Guida, ``Obstacle avoidance system for unmanned ground vehicles by using ultrasonic sensors,'' \textit{Machines}, vol.~6, no.~2, p.~18, 2018.

\bibitem{gronewold2025tactile}
A.~M.~Gronewold, P.~Mulford, E.~Ray, and L.~E.~Ray, ``Tactile sensing \& visually-impaired navigation in densely planted row crops, for precision fertilization by small UGVs,'' \textit{Computers and Electronics in Agriculture}, vol.~231, p.~110003, 2025.

\bibitem{sanchez2023reinforcement}
M.~S\'{a}nchez, J.~Morales, and J.~L.~Mart\'{i}nez, ``Reinforcement and curriculum learning for off-road navigation of an UGV with a 3D LiDAR,'' \textit{Sensors}, vol.~23, no.~6, p.~3239, 2023.

\bibitem{tibermacine2026autonomous}
A.~Tibermacine, I.~E.~Tibermacine, D.~Akrour, A.~Rabehi, and M.~Habib, ``Autonomous navigation in unstructured outdoor environments using semantic segmentation guided reinforcement learning,'' \textit{Scientific Reports}, vol.~16, no.~1, p.~2633, 2026.

\bibitem{mukherjee2025terrain}
J.~Mukherjee, ``Terrain Aware Tactical Motion Planning and Control Algorithms for Off-Road UGVs in GNSS-Denied Hostile Environments,'' Ph.D. dissertation, Virginia Polytechnic Institute and State University, 2025.

\bibitem{jung2026autonomous}
M.~Jung and A.~J.~Choi, ``Autonomous charging system via manipulator-UGV docking using zero-shot 6-DoF pose estimation,'' \textit{Alexandria Engineering Journal}, vol.~134, pp.~122--134, 2026.

\bibitem{naglak2021cable}
J.~E.~Naglak, C.~Kase, M.~McGinty, C.~D.~Majhor, C.~S.~Greene, J.~P.~Bos, and W.~W.~Weaver, ``Cable deployment system for unmanned ground vehicle (UGV) mobile microgrids,'' \textit{HardwareX}, vol.~10, p.~e00205, 2021.

\bibitem{fisher2025developing}
C.~Fisher and M.~J.~McDermott, ``Developing a Communication Node Deployment System for UGVs,'' \textit{IFAC-PapersOnLine}, vol.~59, no.~32, pp.~114--119, 2025.

\bibitem{memet2025adaptive}
D.~Memet, I.~Pirnog, A.-M.~Dr\u{a}gulinescu, G.~Craioveanu, C.~Zamfirescu, and I.~Marcu, ``Adaptive P2P UAV-UGV prototype for ad hoc network-based precision agriculture,'' in \textit{Proc. International Wireless Communications and Mobile Computing (IWCMC)}, 2025, pp.~1552--1557.

\end{thebibliography}


\end{document}
